
\subsection{RQ2: Can existing SATD detection tools detect test code SATD?}
\textbf{Motivation.} As prior research has shown that self-admitted technical debt (SATD) adversely affects software quality and maintenance efforts~\cite{chowdhury_evidence_2025}, the ability to automatically detect SATD is essential for taking timely actions. Yet, it is unknown whether current automated techniques can identify SATD within test code, which is investigated here.

% \textbf{Approach.} We employed six well-established SATD detection approaches originally developed for source code comments. First, the \textit{pattern}-based approach by Potdar et al.~\cite{potdar_exploratory_2014}, which relies on 62 manually curated patterns. Second, the \textit{NLP}-based approach by Maldonado et al.~\cite{maldonado_using_2017}, which applies Natural Language Processing techniques with a maximum entropy classifier. Third, the text mining (\textit{TM}) approach by Huang et al.~\cite{huang_identifying_2018}, which uses NLP preprocessing, feature selection, and a Vector Space Model (VSM). Fourth, the Matches task Annotation Tags (\textit{MAT}) approach by Guo et al.~\cite{guo_how_2021}, which detects SATD by fuzzily matching comment content with common task annotation tags in Integrated Development Environments (IDEs). Finally, we evaluated the BERT-based multitask learning (MT-BERT) approach proposed by Gu et al.~\cite{gu_self-admitted_2024}, which has been shown to outperform existing SATD detection techniques from four different artifacts. Among these approaches, the Pattern and MAT are unsupervised and do not require any training, whereas NLP, TM and MT-BERT are supervised. The pattern-based approach is straightforward to implement. However, since the original authors of the NLP approach did not share their source code, Guo et al.~\cite{guo_how_2021} and Huang et al.~\cite{huang_identifying_2018} re-implemented, and the former showed that their version performed even better. Hence, for reproducibility and fair comparison, we utilized the JAR package provided by Guo et al.~\cite{guo_how_2021}, which encapsulates all implementations except MT-BERT. For MT-BERT, we used the original implementation by Gu et al.~\cite{gu_self-admitted_2024} with the default parameter settings reported in their study. For all implementations, we evaluated the pretrained models on their corresponding datasets and subsequently retrained the models on our test code comment dataset before re-evaluating their performance.

\textbf{Approach.} Seven well-established SATD detection approaches originally developed for source code comments were employed, following the same implementations used in prior comparative work whenever possible.

\textbf{Pattern.} The 62 keyword/phrase patterns summarized by Potdar et al.~\cite{potdar_exploratory_2014} are used. For each target comment, if any predefined pattern is matched, the comment is labeled as SATD; otherwise, it is labeled as non-SATD. This is a strict rule-based unsupervised method and does not require model training.

\textbf{MAT.} Similar to Pattern, Matches task Annotation Tags (MAT)~\cite{guo_how_2021} is unsupervised; however, it differs in two key ways: it is centered on widely used task-annotation tags (e.g., TODO/FIXME-style tags supported by IDEs), and it applies fuzzy matching rather than strict exact matching, enabling more flexible identification of SATD-like expressions in comments.

\textbf{SATD Detector.} The text-mining based SATD Detector by Liu et al.~\cite{liu_satd_2018} is included as a supervised classifier that identifies SATD comments from textual features through preprocessing, feature selection, and vector-space representations.

\textbf{TM.} Another text-mining approach by Huang et al.~\cite{huang_identifying_2018} is also included, using the implementation packaged by Guo et al.~\cite{guo_how_2021}.

\textbf{NLP.} Maldonado et al.~\cite{maldonado_using_2017} formulated SATD detection as an NLP-based text classification task with a Maximum Entropy classifier. As in Guo et al.~\cite{guo_how_2021}, the available reimplementation is used (the original code was not released) to ensure reproducibility and a consistent experimental pipeline.

\textbf{MT-BERT.} For the multitask BERT model proposed by Gu et al.~\cite{gu_self-admitted_2024}, the original released implementation is used with default settings reported by the authors. BERT is a bidirectional transformer-based language model that learns contextual representations from text.

\textbf{DebtHunter. DebtHunter~\cite{sala_debthunter_2021}, a machine-learning-based SATD detection approach that classifies comments using learned textual features, is also evaluated.}

In summary, Pattern, SATD Detector, NLP, TM, and MAT are executed via Guo et al.'s JAR package~\cite{guo_how_2021}, while MT-BERT and DebtHunter are run using their original implementations~\cite{gu_self-admitted_2024,sala_debthunter_2021}. For supervised methods (SATD Detector, NLP, TM, MT-BERT, and DebtHunter), pretrained, retrained, and retrained-with-cross-validation settings are reported where applicable. Precision, Recall, F1-score, and Matthews Correlation Coefficient (MCC) are reported.

\begin{table*}[ht!]
\centering
\caption{SATD detection performance with/without retraining. The unsupervised MAT approach achieved the highest F1-score. Pattern-based methods had high precision but very low recall, and the Text Mining (TM) Method showed low precision but relatively higher recall. We also show the minimum, maximum, and average of all the scores across all the tools and approaches.}
\label{table:non_llm_satd_detection_result}
\begin{tabular}{cccccccc}
\toprule
\textbf{Dataset} & \textbf{Method} & \textbf{Retrained?} & \textbf{Cross-Validation?} & \textbf{Precision} & \textbf{Recall} & \textbf{F1-score} & \textbf{MCC}\\

\midrule
\multirow{17}{*}{\rotatebox[origin=c]{45}{Original}}
  & Pattern~\cite{potdar_exploratory_2014} & -- & --    &   0.80 & 0.03 & 0.06 & 0.15\\
  & MAT~\cite{guo_how_2021} & -- & --    &   \textbf{0.95} & 0.61 & \textbf{0.75} & \textbf{0.76}\\
  & SATD Detector~\cite{liu_satd_2018} & \xmark & \xmark    &   0.38 & 0.58 & 0.46 & 0.46\\
  & SATD Detector~\cite{liu_satd_2018} & \checkmark & \xmark    &   0.30 & 0.57 & 0.40 & 0.41\\
  & SATD Detector~\cite{liu_satd_2018} & \checkmark & \checkmark    &   0.35 & 0.63 & 0.44 & 0.45\\
  & TM~\cite{huang_identifying_2018} & \xmark & \xmark    &   0.09 & 0.56 & 0.15 & 0.20\\
  & TM~\cite{huang_identifying_2018} & \checkmark & \xmark    &   0.10 & 0.86 & 0.17 & 0.26\\
  & TM~\cite{huang_identifying_2018} & \checkmark & \checkmark    &   0.10 & \textbf{0.88} & 0.18 & 0.28\\
  & NLP~\cite{maldonado_using_2017} & \xmark & \xmark    &   0.53 & 0.54 & 0.53 & 0.53\\
  & NLP~\cite{maldonado_using_2017} & \checkmark & \xmark    &   0.68 & 0.63 & 0.65 & 0.65\\
  & NLP~\cite{maldonado_using_2017} & \checkmark & \checkmark    &   0.86 & 0.67 & \textbf{0.75} & 0.75\\
  & MT-BERT~\cite{gu_self-admitted_2024} & \xmark & \xmark    &   0.83 & 0.63 & 0.71 & 0.72\\
  & MT-BERT~\cite{gu_self-admitted_2024} & \checkmark & \xmark    &   0.63 & 0.50 & 0.56 & 0.56\\
  & MT-BERT~\cite{gu_self-admitted_2024} & \checkmark & \checkmark    &   0.68 & 0.45 & 0.54 & 0.55\\
  & DebtHunter~\cite{sala_debthunter_2021} & \xmark & \xmark    &   \added{0.26} & \added{0.58} & \added{0.36} & \added{0.38}\\
  & DebtHunter~\cite{sala_debthunter_2021} & \checkmark & \xmark    &   \added{0.53} & \added{0.65} & \added{0.59} & \added{0.58}\\
  & DebtHunter~\cite{sala_debthunter_2021} & \checkmark & \checkmark    &   \added{0.68} & \added{0.69} & \added{0.68} & \added{0.68}\\
  & \textit{Minimum} & -- & -- & 0.09 & 0.03 & 0.06 & 0.15\\
  & \textit{Maximum} & -- & -- & 0.95 & 0.88 & 0.75 & 0.76\\
  & \textit{Average} & -- & -- & 0.52 & \added{0.59} & \added{0.47} & \added{0.49}\\
\midrule
\multirow{17}{*}{\rotatebox[origin=c]{45}{Deduplicate}}
  & Pattern~\cite{potdar_exploratory_2014} & -- & --    &   0.80 & 0.04 & 0.07 & 0.17\\
  & MAT~\cite{guo_how_2021} & -- & --    &   \textbf{0.95} & 0.56 & \textbf{0.70} & \textbf{0.73}\\
  & SATD Detector~\cite{liu_satd_2018} & \xmark & \xmark    &   0.37 & 0.63 & 0.46 & 0.47\\
  & SATD Detector~\cite{liu_satd_2018} & \checkmark & \xmark    &   0.31 & 0.61 & 0.41 & 0.42\\
  & SATD Detector~\cite{liu_satd_2018} & \checkmark & \checkmark    &   0.36 & \added{0.65} & 0.46 & 0.47\\
  & TM~\cite{huang_identifying_2018} & \xmark & \xmark    &   0.14 & 0.49 & 0.22 & 0.24\\
  & TM~\cite{huang_identifying_2018} & \checkmark & \xmark    &   0.28 & 0.62 & 0.39 & 0.40\\
  & TM~\cite{huang_identifying_2018} & \checkmark & \checkmark    &   0.30 & 0.63 & 0.40 & 0.42\\
  & NLP~\cite{maldonado_using_2017} & \xmark & \xmark    &   0.49 & 0.50 & 0.49 & 0.48\\
  & NLP~\cite{maldonado_using_2017} & \checkmark & \xmark    &   0.64 & 0.58 & 0.61 & 0.60\\
  & NLP~\cite{maldonado_using_2017} & \checkmark & \checkmark    &   0.81 & 0.62 & \textbf{0.70} & 0.70\\
  & MT-BERT~\cite{gu_self-admitted_2024} & \xmark & \xmark    &   0.83 & 0.58 & 0.69 & 0.69\\
  & MT-BERT~\cite{gu_self-admitted_2024} & \checkmark & \xmark    &   0.65 & 0.43 & 0.52 & 0.53\\
  & MT-BERT~\cite{gu_self-admitted_2024} & \checkmark & \checkmark    &   0.60 & 0.41 & 0.49 & 0.49\\
  & DebtHunter~\cite{sala_debthunter_2021} & \xmark & \xmark    &   \added{0.34} & 0\added{.53} & \added{0.41} & \added{0.41}\\
  & DebtHunter~\cite{sala_debthunter_2021} & \checkmark & \xmark    &   \added{0.21} & \added{0.65} & \added{0.32} & \added{0.35}\\
  & DebtHunter~\cite{sala_debthunter_2021} & \checkmark & \checkmark    &   \added{0.47} & \added{\textbf{0.69}} & \added{0.54} & \added{0.55}\\
  & \textit{Minimum} & -- & -- & 0.14 & 0.04 & 0.07 & 0.17\\
  & \textit{Maximum} & -- & -- & 0.95 & \added{0.69} & 0.70 & 0.73\\
  & \textit{Average} & -- & -- & \added{0.50} & \added{0.54} & \added{0.46} & \added{0.48}\\
\bottomrule
\end{tabular}

\end{table*}


\textbf{Results.} The detection results of the seven approaches are summarized in Table~\ref{table:non_llm_satd_detection_result}. The pattern-based approach consistently achieved high precision (0.80) but extremely low recall on both the original (0.03) and deduplicated (0.04) test sets, resulting in very low F1-scores (0.06 and 0.07, respectively). Notably, both precision and F1-score are approximately two to four times lower for the test code dataset than for their original dataset. This outcome highlights the approach’s limited coverage for test code, as it relies on a fixed set of 62 manually curated patterns obtained through a source code-dominated dataset. The low recall can be attributed, in part, to missing keywords in the manually crafted list, specifically the absence of \texttt{TODO} alongside \texttt{FIX}, as well as the lack of discernible patterns in many comments.

The MAT approach achieved the best overall performance across both datasets without any training. It consistently yielded high precision (0.95) and achieved the highest F1-scores on both the original (0.75) and deduplicated (0.70) datasets, with corresponding recall values of 0.61 and 0.56. Notably, both MAT and the NLP approach achieved the highest F1-scores, with NLP attaining this performance under the cross-validation setting. Nevertheless, compared to prior work, the F1-score observed in this evaluation is 2\% lower.

The two text-mining approaches exhibit notably different performance profiles on test code comments. Compared with Huang et al.’s TM approach, SATD Detector consistently achieves much higher precision, F1-score, and MCC on both datasets, while maintaining competitive recall. However, SATD Detector remains clearly below the retrained NLP approach, which attains substantially stronger balanced performance, indicating that SATD Detector provides better coverage-oriented performance than TM but not the same overall effectiveness as NLP. Meanwhile, the TM approach produced consistently low precision and F1-scores across all settings. Similar to the NLP pretrained model, TM failed to recognize \texttt{TODO}, although it correctly identified \texttt{FIXME}, and produced numerous false positives from patterns such as issue IDs (e.g., \texttt{// DROOLS-155}), copyright notices, and Javadocs with author annotations. After retraining, TM corrected many of these patterns but still misclassified a substantial number of irrelevant comments as SATD. Nevertheless, its recall was comparable to that of the NLP and MAT approaches, and was even the highest for the retrained cross-validation setting on the original dataset. Compared to Huang et al.’s~\cite{huang_identifying_2018} reported F1-score of 0.737 for source code, its performance on test code is notably poor.

For both datasets, the retrained NLP model achieved roughly 12\% higher F1-scores than the pretrained model, producing balanced precision and recall. On the pretrained model, many comments containing the two major keywords, \texttt{TODO} and \texttt{FIXME}, were missed, and the misidentified comments generally lacked clear patterns. After retraining, the model successfully learned \texttt{TODO} but still failed to capture \texttt{FIXME}. Notably, most of the remaining misclassified comments were Javadocs describing bug-related tests that include bug ID (e.g., \texttt{@bug 8077931}). Compared to Guo et al.’s~\cite{guo_how_2021} result on source code, the best F1-score in this evaluation after retraining is about 5\% lower.

The MT-BERT approach exhibited strong performance on pretrained, achieving F1-scores of 0.71 and 0.69 on the original and deduplicated test datasets, respectively. This indicates that the pretrained multitask model generalizes reasonably well to test code comments. However, retraining MT-BERT on the test code comment datasets resulted in a noticeable performance degradation, primarily due to reduced recall, which led to lower F1-scores. Manual inspection revealed that MT-BERT attained relatively high precision by correctly identifying SATD comments containing explicit debt-related keywords; nevertheless, it occasionally failed to detect when common keywords such as \texttt{TODO} were present, and similar to MAT, it struggled with comments lacking explicit keyword signals, leading to comparable recall between the two approaches. Overall, the pretrained version performed better in these cases, likely because it was trained on a larger and more diverse set of SATD patterns, enabling it to better capture keyword-based and recurring SATD expressions. DebtHunter improved after retraining, achieving stronger recall and F1-score under cross-validation. However, its lower precision kept its overall performance below the best MAT and NLP results.

The overall poor to moderate performance indicates that methods originally developed based on the characteristics of source code fail to fully capture the distinctive properties of test code. This finding reinforces the notion that SATD instances in test code differ substantially from those in source code. Although retraining sometimes improved the performance, even in the best scenarios, results remained 2–5\% lower than those on the source code. These findings highlight the somewhat poor transferability of source-trained models and show that retraining alone is not sufficient.

\begin{tcolorbox}
\noindent\textbf{\underline{RQ2 summary}:}
Although the NLP and MAT approaches achieve the highest accuracy in detecting SATD in test code, their F1-scores remain relatively modest at around 0.75. This limitation motivated the investigation of different LLMs for detecting SATD in test code, particularly because prior research has demonstrated the strong performance of LLMs across various software engineering tasks, including the detection of SATD in source code.
\end{tcolorbox}
